\begin{theorem}[Iterator Successor Clause]
\label{thm:iterator-successor-step}
Let $(P,S,1)$ be a Peano system, let $(W,c,g)$ be iterator data for it, and let
$f$ be the iterator-generated function determined by $(W,c,g)$. Then
\[
\forall n\in P\;(f(S(n))=g(f(n))).
\]
\hyperref[prf:iterator-successor-step]{Go to proof.}
\end{theorem}
\end{theorembox}

\begin{remark*}[Standard quantified statement]
\[
f \text{ is the iterator-generated function determined by } (W,c,g)
\Longrightarrow
\forall n\in P\;(f(S(n))=g(f(n))).
\]
\end{remark*}

\begin{remark*}[Predicate reading]
\[
\operatorname{IteratorGeneratedFunction}(f,(P,S,1),W,c,g)
\Longrightarrow
\forall n\in P\;(f(S(n))=g(f(n))).
\]
\end{remark*}

\begin{remark*}[Contrapositive quantified statement]
\[
\exists n\in P\;(f(S(n))\neq g(f(n)))
\Longrightarrow
f \text{ is not the iterator-generated function determined by } (W,c,g).
\]
\end{remark*}

\begin{remark*}[Interpretation]
The successor clause says that the generated function advances by applying
the update rule to the previous value at every Peano stage.
\end{remark*}

\begin{dependencies}
\begin{itemize}
  \item \hyperref[def:iterator-generated-function]{Iterator-Generated Function}
\end{itemize}
\end{dependencies}

